\chapter{Nhật ký thực nghiệm và các lần chạy thất bại}
\label{app:nhat-ky}

Phụ lục này ghi lại các lần chạy chính, thời lượng, trạng thái và bài học từ các lần thất bại.
Mục đích không phải kể lể, mà để người đọc biết con số nào đến từ run hoàn chỉnh, con số nào
đến từ run bị ngắt.

\section{Dòng thời gian}
\begin{center}\small
\begin{longtable}{@{}l l r r@{}}
\toprule
\textbf{Thời điểm} & \textbf{Run} & \textbf{Thời lượng} & \textbf{Ván} \\
\midrule
14/11/2026 & A3 MCTS số mô phỏng & 22 phút & 1\,600 \\
14/11/2026 & A4 rollout & 9 phút & 800 \\
17/11/2026 & RL smoke 1 & 28 phút (ngắt) & 0 \\
17/11/2026 & RL smoke 2 & 4 phút & 0 \\
17/11/2026 & RL smoke 3 & 4 phút & 0 \\
17/11/2026 & RL v2 (từ SL v2) & 47 phút & 49\,920 episode \\
18/11/2026 & RL v3 (từ SL v3) & 46 phút & 49\,920 episode \\
23/11/2026 & L8 smoke & 2 giây & 6 \\
23/11/2026 & L8 GPU (ngắt) & 2 giờ 44 phút & 3\,007 \\
24--25/11/2026 & L8 GPU native & 12 giờ 19 phút & 15\,020 \\
25/11/2026 & L8 gặp L3 & 2,6 phút & 40 \\
25/11/2026 & E3 tự đấu & 7,3 phút & 800 \\
\bottomrule
\end{longtable}
\end{center}

\section{Các lần chạy thất bại}
\paragraph{RL smoke 1 bị ngắt.} Run \texttt{20261117-164022\_smoke} bị ngắt sau 28 phút. Nguyên
nhân là cấu hình thử quá lớn so với thời gian chờ; không có kết quả dùng trong báo cáo. Bài học:
chạy smoke với vài ván trước khi chạy 50k episode.

\paragraph{L8 GPU bị ngắt.} Run \texttt{20261123-194550\_gpu} chạy 2 giờ 44 phút, được 3\,007
ván thì bị ngắt (KeyboardInterrupt trong \texttt{\_expand}). Nguyên nhân có thể là người dùng
dừng để đổi cấu hình backend. Run này có \texttt{metrics.jsonl} nhưng không có
\texttt{summary.json} trạng thái ok; báo cáo không dùng làm bằng chứng chính. Run kế tiếp
\texttt{20261124-160615\_gpu-native} chạy đủ 50 vòng.

\paragraph{Lỗi quy ước value.} Không phải run thất bại mà là lỗi logic: value loss của RL luôn
1,2--1,4, tệ hơn đoán 0. Đã phân tích, sửa ở SL v3, và đo lại. Đây là bài học lớn nhất của dự án.

\paragraph{Đường dẫn input sai trong báo cáo.} Bản nháp đầu dùng \texttt{\textbackslash
input\{so-lieu-...\}} thay vì \texttt{ban-mo-rong/so-lieu-...}, làm build lỗi. Đã sửa. Bài học:
build sớm, đừng để đến cuối.

\paragraph{Tham số model lệch YAML.} File YAML ghi \texttt{fc\_hidden: 64}, code dùng 32. Không
ảnh hưởng kết quả nhưng gây nhầm khi đọc. Đã ghi chú.

\section{Chi phí}
\begin{center}\small
\begin{tabularx}{\textwidth}{@{}l r X@{}}
\toprule
\textbf{Hạng mục} & \textbf{Thời gian} & \textbf{Ghi chú} \\
\midrule
Sinh dữ liệu chuyên gia v2 & vài giờ & 297\,098 mẫu, 20\,000 ván \\
SL v3 & khoảng 20--30 phút & 60 epoch, GPU \\
RL v3 & 46 phút & 49\,920 episode \\
L8 & 12 giờ 19 phút & 10\,000 ván tự đấu + 50 gating \\
RR8 & khoảng 30 phút & 2\,800 ván, 1 worker \\
E3 & 7,3 phút & 800 ván \\
A3 & 22 phút & 1\,600 ván \\
A4 & 9 phút & 800 ván \\
\bottomrule
\end{tabularx}
\end{center}
Chi phí lớn nhất là L8, chủ yếu ở tự đấu (67\%) và gating (19\%). Nếu có GPU mạnh hơn, có thể
tăng số ván hoặc số mô phỏng, nhưng kết quả sẽ khác; báo cáo chỉ ghi cấu hình đã chạy.

\section{Bài học vận hành}
\begin{enumerate}
  \item \textbf{Chạy smoke trước.} Mọi cấu hình mới nên chạy 2--6 ván để kiểm tra pipeline,
        checkpoint và log trước khi chạy hàng nghìn ván.
  \item \textbf{Ghim checkpoint.} Không dùng \texttt{latest} trong báo cáo; luôn dùng đường dẫn
        có timestamp hoặc revision.
  \item \textbf{Ghi backend.} Cùng thuật toán nhưng Python/Cython/GPU cho thời gian khác; không
        so thời gian giữa các run khác backend.
  \item \textbf{Giữ log dirty.} Run chạy khi working tree dirty vẫn có ích để gỡ lỗi, nhưng
        không dùng làm bằng chứng chính. \texttt{meta.json} giúp phân biệt.
  \item \textbf{Đừng chạy nặng khi người dùng cần máy.} Các run lớn nên chạy với \texttt{nice},
        \texttt{ionice} và giới hạn CPU; hoặc chạy qua đêm.
\end{enumerate}